% report.tex
\documentclass[dvipdfmx, fontsize=12pt]{jlreq}
%\documentclass[dvipdfmx, a4paper, 14pt]{jsarticle}

\usepackage{graphicx}
\usepackage{tikz}
%\usepackage{bxpgfcurpage}
\usepackage{lipsum}
%\usepackage{xsavebox} %save content for repeated use
%\usepackage{atbegshi} %insert material on every page
\usepackage{layout}
\usepackage[top=2.3truecm, bottom=1.3truecm, left=0.6truecm, right=0.6truecm, footskip=0.8truecm, headsep=1.5truecm]{geometry}
\usetikzlibrary{calc}
%\usepackage[T1]{fontenc}
\usepackage{newpxtext, newpxmath}
%\usepackage{lmodern}
\usepackage{fancyhdr}
\usepackage{lastpage}
\usepackage{eso-pic}
%\usepackage[hiragino-pron]{pxchfon}


\renewcommand\kanjifamilydefault{\gtdefault}
\renewcommand\familydefault{\sfdefault}

\pagestyle{fancy}
% ヘッダ右
\renewcommand{\headrulewidth}{0pt} % ヘッダーの羅線をなくす
\fancyhead[R]{\footnotesize 全ページ \pageref{LastPage} 中 \thepage{} ページ}
% フッタ中央
\fancyfoot[C]{\small hogehoge\&foobar株式会社}

%%%%%%%%%%%%%%
% 報告書No.はここを変える
\newcommand{\登録番号}{TR23-1001094-01}
%%%%%%%%%%%%%%

% 画像フォルダのパスを設定
\graphicspath{{./image/}}


%% 参考：https://dayinthelife-web.blogspot.com/2014/01/latex-geometry.html
\newsavebox{\myheader}
\savebox{\myheader}{%
\begin{tikzpicture}[overlay]
\draw [line width=1.0pt]
($ (current page.north west) + (0.3cm,-1.0cm) $)
rectangle
($ (current page.north west) + (12.0cm,-2.0cm) $);
\draw ($ (current page.north west) + (6.5cm,-1.5cm) $) node{\large \textbf{技 術 報 告 書}};
\draw [line width=1.0pt]
($ (current page.north west) + (12.0cm,-1.0cm) $)
rectangle
($ (current page.north west) + (15.0cm,-2.0cm) $);
\draw ($ (current page.north west) + (13.5cm,-1.5cm) $) node{\textbf{登録番号}};
\draw [line width=1.0pt]
($ (current page.north west) + (15.0cm,-1.0cm) $)
rectangle
($ (current page.north east) + (-0.3cm,-2.0cm) $);
\draw ($ (current page.north east) + (-3.2cm,-1.5cm) $) node{\textsf{\登録番号}};

%% 本文部分の枠組み
\draw [line width=1.0pt]
($ (current page.north west) + (0.3cm,-2.0cm) $)
rectangle
($ (current page.south east) + (-0.3cm,1.0cm) $);
\end{tikzpicture}%
}

\AddToShipoutPictureBG{\usebox{\myheader}}


% ここから本文
\begin{document}

\begin{center}
        \vspace*{1cm}

        \huge
        \textbf{Sample title page with border}

        \vspace{2.0cm}
        \LARGE
        \textbf{First name, Family name}

        \vfill
\end{center}
%\end{titlepage}
\newpage

\begin{tikzpicture}[remember picture, overlay]
    \node[color=red!30!white,draw,line width=10pt,
    rectangle,font=\gtfamily,scale=10,rotate=30]
     at (current page.center) {部外秘};
\end{tikzpicture}
\lipsum[1-5]

\newpage

\section{はじめに}
これは報告書用のフォーマットです。

\subsection{使い方}
\begin{itemize}
    \item これは報告書用のフォーマットです
    \item これは報告書用のフォーマットです
\end{itemize}

図表はこんな感じになります。図\ref{tiger}にトラの絵を示した。

45PL $\times$ 40Aの組み合わせにて検証を行った。
\begin{figure}
    \centering
    \includegraphics[keepaspectratio=true,width=0.47\textwidth]{tiger.png}
    \caption{This is a tiger.}
    \label{tiger}
\end{figure}

\end{document}
